\documentclass[a4paper,11pt,oneside]{book}
% preamble_build.tex: copy of Maximus's preamble.tex for the DE v3 build.
% The blocks marked [BUILD COPY CHANGE] are the earlier build fixes.
% The blocks marked [V3 ADDITION] are new in this version (example box,
% spacing). Everything else is unchanged from preamble.tex.

% For margins
\usepackage[a4paper, margin = 20mm]{geometry}

\usepackage[Lenny]{fncychap}
% Rejne, Glenn, Lenny

\usepackage{etoc}

% [V3 ADDITION] cancel: used to show the dx's "cancelling" in Section 11.3
\usepackage{cancel}
\renewcommand{\CancelColor}{\color{red}}

% For enumeration & math symbols
\usepackage{enumerate}
\usepackage{amsmath,amsfonts,amssymb,amsthm,graphicx,lipsum}

% For diagrams, plots and boxes
\usepackage{pgfplots}
\pgfplotsset{compat=1.18}
\usepackage{tikz}
\usepackage[most]{tcolorbox}
\tcbuselibrary{skins}
\usepackage{xcolor}
\pgfplotsset{
    standard/.style={clip = false,
    axis line style = thick,
    trig format=rad,
    axis lines = middle,
    every axis x label/.style={at={(current axis.right of origin)},anchor= west},
    every axis y label/.style={at={(current axis.above origin)},anchor=south}
    }
}

% For GC font
\usepackage{fontspec}
% [BUILD COPY CHANGE 1] The TI-84 key font is loaded only if the file exists
% (checked in ./Fonts/... and ../Fonts/...). Otherwise \ti{} falls back to a
% small boxed sans-serif key. Known letters of the TI font are mapped to their
% key names in the fallback: o = [y=], z = [mode], f = [L3].
\newcommand{\tikey}[1]{\tcbox[on line, enhanced, boxsep=0pt, left=2pt, right=2pt,
  top=1pt, bottom=1pt, boxrule=0.5pt, arc=1.5pt, colback=white,
  colframe=black!70]{\sffamily\footnotesize\strut #1}}
\ExplSyntaxOn
\NewDocumentCommand{\tifallback}{m}{
  \str_case:nnF {#1}
    { {o}{\tikey{y=}} {z}{\tikey{mode}} {f}{\tikey{L3}} }
    { \tikey{#1} }
}
\ExplSyntaxOff
\def\tifontcase{0}
\IfFileExists{Fonts/TI-84 Plus CE Key Fonts/TI84PlusCEKeys.ttf}
  {\def\tifontcase{1}}{}
\IfFileExists{../Fonts/TI-84 Plus CE Key Fonts/TI84PlusCEKeys.ttf}
  {\def\tifontcase{2}}{}
\ifcase\tifontcase\relax
  % 0: font not found, use boxed fallback keys
  \newcommand{\ti}[1]{\tifallback{#1}}
\or
  % 1: original location (as in preamble.tex)
  \newfontfamily\tifont[
    Path = {Fonts/TI-84 Plus CE Key Fonts/},
    UprightFont = {TI84PlusCEKeys.ttf}
  ]{}
  \newcommand{\ti}[1]{{\tifont #1}}
\or
  % 2: one level up (when compiling inside draft-ch1/)
  \newfontfamily\tifont[
    Path = {../Fonts/TI-84 Plus CE Key Fonts/},
    UprightFont = {TI84PlusCEKeys.ttf}
  ]{}
  \newcommand{\ti}[1]{{\tifont #1}}
\fi

% SI units
\usepackage{siunitx}

% setspace
%\usepackage{setspace}
%\onehalfspacing

% For hyperlinks within and out of document
\usepackage{hyperref}
\hypersetup{
    colorlinks = true,
    linkcolor = black,
    urlcolor = blue,
}   

% For captions for pics
\usepackage{capt-of}


% tcolorbox custom environments
% Definition ========================================================
\newtcolorbox[auto counter,number within=section]{definition}[2][drop shadow]{%
enhanced, before skip=5mm, after skip=6mm, % [V3 ADDITION] spacing
colback=blue!5!white,colframe=blue!50!black,fonttitle=\bfseries,
title= Definition~\thetcbcounter. #2,#1}

% Theorem ==========================================================
\newtcolorbox[auto counter,number within=section]{theorem}[2][drop shadow]{%
enhanced, before skip=5mm, after skip=6mm, % [V3 ADDITION] spacing
colback=black!5!white,colframe=black,fonttitle=\bfseries,
title= Theorem~\thetcbcounter. #2,#1}

% Caution box ======================================================
\newtcolorbox{caution}[1][]{enhanced,
before skip=5mm,after skip=6mm, % [V3 ADDITION] was 2mm/3mm
boxrule=1pt,left=11mm,right=2mm,top=3mm,bottom=3mm,
colback=yellow!50,
colframe=yellow!20!black,
sharp corners,rounded corners=southeast,arc is angular,arc=3mm,
underlay={%
\path[fill=tcbcolback!80!black] ([yshift=3mm]interior.south east)--++(-0.4,-0.1)--++(0.1,-0.2);
\path[draw=tcbcolframe,shorten <=-0.05mm,shorten >=-0.05mm] ([yshift=3mm]interior.south east)--++(-0.4,-0.1)--++(0.1,-0.2);
\path[fill=yellow!50!black,draw=none] (interior.south west) rectangle node[white]{\Huge\bfseries !} ([xshift=9mm]interior.north west);
},
drop fuzzy shadow,#1}

% Question box ======================================================
\newtcolorbox{question}[1][]{enhanced,
before skip=5mm,after skip=6mm, % [V3 ADDITION] was 2mm/3mm
boxrule=1pt,left=11mm,right=2mm,top=3mm,bottom=3mm,
colback=blue!20!white,
colframe=blue!20!black,
sharp corners,rounded corners=southeast,arc is angular,arc=3mm,
underlay={%
\path[fill=tcbcolback!80!black] ([yshift=3mm]interior.south east)--++(-0.4,-0.1)--++(0.1,-0.2);
\path[draw=tcbcolframe,shorten <=-0.05mm,shorten >=-0.05mm] ([yshift=3mm]interior.south east)--++(-0.4,-0.1)--++(0.1,-0.2);
\path[fill=blue!50!black,draw=none] (interior.south west) rectangle node[white]{\Huge\bfseries ?} ([xshift=9mm]interior.north west);
},
drop fuzzy shadow,#1}

% Information box ======================================================
\newtcolorbox{info}[1][]{enhanced,
before skip=5mm,after skip=6mm, % [V3 ADDITION] was 2mm/3mm
boxrule=1pt,left=11mm,right=2mm,top=3mm,bottom=3mm,
colback=cyan!20!white,
colframe=cyan!20!black,
sharp corners,rounded corners=southeast,arc is angular,arc=3mm,
underlay={%
\path[fill=tcbcolback!80!black] ([yshift=3mm]interior.south east)--++(-0.4,-0.1)--++(0.1,-0.2);
\path[draw=tcbcolframe,shorten <=-0.05mm,shorten >=-0.05mm] ([yshift=3mm]interior.south east)--++(-0.4,-0.1)--++(0.1,-0.2);
\path[fill=cyan!50!black,draw=none] (interior.south west) rectangle node[white]{\Huge\bfseries i} ([xshift=9mm]interior.north west);
},
drop fuzzy shadow,#1}

% Mistake box ======================================================
\newtcolorbox{mistake}[1][]{enhanced,
before skip=5mm,after skip=6mm, % [V3 ADDITION] was 2mm/3mm
boxrule=1pt,left=11mm,right=2mm,top=3mm,bottom=3mm,
colback=red!20!white,
colframe=red!20!black,
sharp corners,rounded corners=southeast,arc is angular,arc=3mm,
underlay={%
\path[fill=tcbcolback!80!black] ([yshift=3mm]interior.south east)--++(-0.4,-0.1)--++(0.1,-0.2);
\path[draw=tcbcolframe,shorten <=-0.05mm,shorten >=-0.05mm] ([yshift=3mm]interior.south east)--++(-0.4,-0.1)--++(0.1,-0.2);
\path[fill=red!50!black,draw=none] (interior.south west) rectangle node[white]{\Huge\bfseries $\times$} ([xshift=9mm]interior.north west);
},
drop fuzzy shadow,#1}



% [V3 ADDITION] Worked example box ==================================
% Own counter, numbered within the section (Example 11.2.1, 11.2.2, ...).
% For per-subsection numbering (11.2.1.1) change `section' to `subsection'. Blue `question' boxes are kept for thought-provoking questions
% only; worked examples always use this box. Usage:
%   \begin{example} ...question text... \end{example}
%   \solution{...}
\newtcolorbox[auto counter,number within=section]{example}[1][]{%
enhanced, breakable=false, before skip=5mm, after skip=4mm,
colback=teal!5!white, colframe=teal!55!black, boxrule=0.8pt,
left=3mm, right=3mm, top=2mm, bottom=2mm,
sharp corners, rounded corners=northwest, arc=3mm,
fonttitle=\bfseries, coltitle=white,
title=Example~\thetcbcounter, attach boxed title to top left={yshift=-2mm, xshift=3mm},
boxed title style={colback=teal!55!black, sharp corners, rounded corners=northwest, arc=2mm},
drop fuzzy shadow, #1}

% [V3 ADDITION] Spacing between sections, subsections, parts, paragraphs =
% (fncychap's chapter heading is untouched; only \section and below.)
\makeatletter
\renewcommand\section{\@startsection{section}{1}{\z@}%
  {-2ex \@plus -.5ex \@minus -.2ex}{3.2ex \@plus .3ex}%
  {\normalfont\LARGE\bfseries\color{blue!50!black}}}
\renewcommand\subsection{\@startsection{subsection}{2}{\z@}%
  {-1.5ex \@plus -.5ex \@minus -.2ex}{2.6ex \@plus .3ex}%
  {\normalfont\Large\bfseries}}
\renewcommand\subsubsection{\@startsection{subsubsection}{3}{\z@}%
  {-4ex \@plus -1ex \@minus -.5ex}{1.6ex \@plus .2ex}%
  {\normalfont\large\bfseries\color{teal!50!black}}}
\makeatother
\setlength{\parindent}{0pt}
\setlength{\parskip}{0.9ex plus 0.3ex minus 0.2ex}

% Headers and Footers
\usepackage{fancyhdr,lastpage}
% [V3 R2] grey left-margin line on the out-of-syllabus pages (item 13)
\usepackage{eso-pic}
\newcommand{\oosmargin}{\AddToShipoutPictureBG{\put(12mm,0){\color{gray!75}\rule{1.6pt}{\paperheight}}}}
\pagestyle{fancy}
\lhead{Maximus Pung Jun}
% [BUILD COPY CHANGE 2] Header bug fix: \chaptername is the word "Chapter",
% so the original header printed "Chapter 1: Chapter". \chaptermark now stores
% the bare chapter title in \leftmark, and the header uses \leftmark.
\renewcommand{\chaptermark}[1]{\markboth{#1}{}}
\rhead{Chapter \thechapter: \leftmark}
% \cfoot{Page \thepage\  of \pageref{LastPage}} % buggy, does not appear on chapter pages.

% [V3 R2] tick and cross markers for mistake boxes (G3)
\newcommand{\wrong}{\textcolor{red!70!black}{\ensuremath{\boldsymbol{\times}}}}
\newcommand{\cmark}{\textcolor{green!50!black}{\ensuremath{\checkmark}}}

\newcommand{\xmark}{\wrong}   % alias: red cross
% [V3 R2] out-of-syllabus helpers (DE notes item 13)
%   \oosmargin    : grey line in the left margin on every following page
%   \oosmarginoff : stop the margin line
%   \oosnote      : grey banner "Out of syllabus. For your interest only."
%   outofsyllabus : environment = new page + margin line + banner; ends with
%                   \clearpage and switches the margin line off again
\newcommand{\oosmarginoff}{\ClearShipoutPictureBG}
\newcommand{\oosnote}[1][Nothing in this section is examined in H2 Mathematics.]{\par\noindent\colorbox{gray!15}{\parbox{\dimexpr\textwidth-2\fboxsep}{\textbf{Out of syllabus. For your interest only.} #1}}\par}
\newenvironment{outofsyllabus}{\clearpage\oosmargin}{\clearpage\oosmarginoff}

% Notation as in Maximus's notes: italic d, display style everywhere (G1).
% Fractions inside exponents stay small, so never use these macros inside ^{}.
\newcommand{\dydx}{\dfrac{dy}{dx}}
\newcommand{\dvdx}{\dfrac{dv}{dx}}
\newcommand{\dudx}{\dfrac{du}{dx}}
\newcommand{\ddydx}{\dfrac{d^{2}y}{dx^{2}}}
\newcommand{\ddt}[1]{\dfrac{d#1}{dt}}
% Small syllabus flag (no longer used in the DE notes, kept for other chapters)
\newcommand{\syllabusflag}[1]{\par\noindent{\footnotesize\color{red!70!black}\textbf{Syllabus note.} #1}\par}


% Teacher's copy toggle ================================================
% [BUILD COPY CHANGE 3] \ifx compares the whole word (the original
% \if\conform test was always false). DE_notes_student.tex predefines
% \conform as HideSolutions.
% A student copy defines \conform as HideSolutions BEFORE loading this preamble.
\providecommand\conform{ShowSolutions}
\newcommand\ShowSolutionsKey{ShowSolutions}
% [V3 ADDITION] In the student copy the hidden solution leaves stretchable
% blank space (\vfill), so the free space on each page is shared out as
% working space below the examples instead of collecting at the bottom.
\newcommand\solution[1]{\ifx\conform\ShowSolutionsKey
  \par\noindent\textbf{Solution.} #1\par\medskip\else\par\vspace{\workspace}\fi}
% [V3 R2] minimum working space left under each example in the student copy
\newlength\workspace\setlength\workspace{3cm plus 1fill}


\newcommand{\NN}{\mathbb{N}}
\newcommand{\ZZ}{\mathbb{Z}}
\newcommand{\QQ}{\mathbb{Q}}
\newcommand{\RR}{\mathbb{R}}
\newcommand{\CC}{\mathbb{C}}


% =====================================================================
\usepackage{pifont}
% preamble_local.tex: extra packages and styles for the Vectors chapter
% ONLY. Loaded after preamble_build.tex (a verbatim copy of the shared
% book preamble from the PnC / DE build, so that it can be re-synced by
% copying it over again).
% =====================================================================
\usepackage{array, booktabs}
\usetikzlibrary{arrows.meta, calc, positioning, 3d, decorations.pathreplacing, angles, quotes}
\input{../vec_common}% vecfig.tex: shared 3D diagram style (tikz-3dplot) for Vectors I / II notes and slides.
\usepackage{tikz-3dplot}
\tdplotsetmaincoords{62}{55}
\definecolor{vA}{RGB}{31,90,180}   % first vector / line 1
\definecolor{vB}{RGB}{210,90,20}   % second vector / line 2
\definecolor{vC}{RGB}{0,128,110}   % result / third
\definecolor{vN}{RGB}{140,40,150}  % normals
\definecolor{vP}{RGB}{200,215,240} % plane 1
\definecolor{vQ}{RGB}{245,215,190} % plane 2
\definecolor{vR}{RGB}{200,235,215} % plane 3
\tikzset{
  v3/.style={tdplot_main_coords, line cap=round, line join=round, >=Stealth, scale=#1},
  v3/.default=1,
  ax/.style={->, black!50, thin},
  va/.style={->, very thick, vA},
  vb/.style={->, very thick, vB},
  vc/.style={->, very thick, vC},
  vn/.style={->, very thick, vN},
  ln/.style={very thick},
  dsh/.style={dashed, black!55},
  pp/.style={fill=vP, fill opacity=0.55, draw=vA!70!black, thin},
  pq/.style={fill=vQ, fill opacity=0.55, draw=vB!70!black, thin},
  pr/.style={fill=vR, fill opacity=0.55, draw=vC!70!black, thin},
  dot/.style={circle, fill=black, inner sep=1.2pt},
}
% \vaxes{xlen}{ylen}{zlen}
\newcommand{\vaxes}[3]{%
  \draw[ax] (0,0,0) -- (#1,0,0) node[anchor=west]{$x$};
  \draw[ax] (0,0,0) -- (0,#2,0) node[anchor=south west]{$y$};
  \draw[ax] (0,0,0) -- (0,0,#3) node[anchor=south]{$z$};
  \node[anchor=east] at (0,0,0) {\scriptsize $O$};}
% \rightmark{corner}{along1}{along2} small right-angle square
\newcommand{\rtang}[3]{\draw[thin] ($(#1)!0.25cm!(#2)$) -- ($($(#1)!0.25cm!(#2)$)+($(#1)!0.25cm!(#3)$)-(#1)$) -- ($(#1)!0.25cm!(#3)$);}
\newenvironment{vis}{\par\medskip\begin{center}}{\end{center}\medskip}
% \vang{vertex}{p1}{p2}{label}: small angle arc drawn on the page
\newcommand{\vang}[4]{%
  \coordinate (vg1) at ($(#1)!0.55cm!(#2)$); \coordinate (vg2) at ($(#1)!0.55cm!(#3)$);
  \coordinate (vgm) at ($(vg1)!0.5!(vg2)$);
  \draw[vC, thick] (vg1) .. controls ($(#1)!0.72cm!(vgm)$) .. (vg2);
  \node[vC] at ($(#1)!0.95cm!(vgm)$) {#4};}


% A little air between the parts of an example, and between lines of a list.
\setlength{\itemsep}{0.85ex}
\setlength{\parsep}{0.35ex}

% Student copy: minimum working space under each example. 2.2cm (instead of
% the book's 3cm) is meant to keep the student copy on the same page breaks
% as the teacher copy.
\setlength\workspace{2.2cm plus 1fill}

\usepackage{needspace,longtable}
\raggedbottom
\setlength{\parindent}{0pt}
\rhead{Vectors I: Basics Worksheet}
\cfoot{\thepage}
\renewcommand{\headrulewidth}{0.4pt}
\newcounter{wq}[section]
\newcommand{\Q}[1]{\par\needspace{5\baselineskip}\medskip\stepcounter{wq}\noindent\textbf{\thewq}\quad #1\par}
\newcommand{\wsec}[1]{\par\needspace{8\baselineskip}\bigskip\stepcounter{section}\noindent{\Large\scshape \thesection\quad #1}\par\medskip}
\renewcommand{\thesection}{\arabic{section}}
\begin{document}
\thispagestyle{fancy}
{\LARGE\scshape Vectors I Basics Worksheet: Questions}\par\smallskip
{Compiled by Maximus Pung Jun}\par\smallskip
{\small\textcolor{black!60}{Basics of vectors: $\mathbf i$, $\mathbf j$, $\mathbf k$ notation, magnitude, unit vectors, parallel vectors, collinear points and the ratio theorem. Each section has footing questions, one tricky question and a True/False check.}}\par\medskip
\wsec{$\mathbf{i}$, $\mathbf{j}$, $\mathbf{k}$ Notation and Column Vectors}
\Q{Write $\mathbf a=3\mathbf i-2\mathbf j+5\mathbf k$ as a column vector, and write $\mathbf b=\col{4}{0}{-1}$ in terms of $\mathbf i$, $\mathbf j$ and $\mathbf k$.}
\vspace{2.6cm}
\Q{The points $A$ and $B$ have coordinates $(1,2,3)$ and $(4,6,1)$. Find $\overrightarrow{AB}$ in terms of $\mathbf i$, $\mathbf j$ and $\mathbf k$.}
\begin{center}\begin{tikzpicture}[x={(0.8cm,0cm)},y={(0.4cm,0.32cm)},z={(0cm,0.8cm)},>=Stealth,line cap=round]
\draw[ax](0,0,0)--(5.8,0,0) node[right]{\scriptsize$x$};
\draw[ax](0,0,0)--(0,7,0) node[above right]{\scriptsize$y$};
\draw[ax](0,0,0)--(0,0,3.8) node[above]{\scriptsize$z$};
\node[below left] at (0,0,0){\scriptsize$O$};
\draw[dsh](1,2,3)--(1,2,0)--(1,0,0);\draw[dsh](1,2,0)--(0,2,0);
\draw[dsh](4,6,1)--(4,6,0)--(4,0,0);\draw[dsh](4,6,0)--(0,6,0);
\draw[va](0,0,0)--(1,2,3);\draw[vb](0,0,0)--(4,6,1);\draw[vc](1,2,3)--(4,6,1);
\node[dot,label=above:{\scriptsize$A(1,2,3)$}] at (1,2,3){};
\node[dot,label=right:{\scriptsize$B(4,6,1)$}] at (4,6,1){};
\end{tikzpicture}\end{center}
\vspace{2.6cm}
\Q{\textbf{(Tricky)} Given $\mathbf p=2\mathbf i+\mathbf k$ and $\mathbf q=-\mathbf i+3\mathbf j$, find $\mathbf r$ such that $2\mathbf p-3\mathbf q+\mathbf r=\mathbf 0$. Give $\mathbf r$ in terms of $\mathbf i$, $\mathbf j$ and $\mathbf k$.}
\vspace{2.6cm}
\par\needspace{8\baselineskip}\medskip\noindent\textbf{Check Your Understanding.} Circle T (True) or F (False) for each statement.\par\smallskip{\renewcommand{\arraystretch}{1.5}\noindent\makebox[\textwidth][c]{\begin{tabular}{|c|>{\centering\arraybackslash}p{12cm}|c|}\hline
 & \textbf{Statement} & \textbf{T/F}\\\hline
(a) & $3\mathbf i+2\mathbf k=\col{3}{2}{0}$ & T\,/\,F \\[1.2ex]\hline
(b) & The $\mathbf j$-component of $4\mathbf i-\mathbf j+2\mathbf k$ is $-1$ & T\,/\,F \\[1.2ex]\hline
(c) & The position vector of $P(1,-2,0)$ is $\mathbf i-2\mathbf j$ & T\,/\,F \\[1.2ex]\hline
(d) & $\overrightarrow{AB}=\overrightarrow{OA}-\overrightarrow{OB}$ & T\,/\,F \\[1.2ex]\hline
(e) & $\mathbf i+\mathbf j+\mathbf k=\col{1}{1}{1}$ & T\,/\,F \\[1.2ex]\hline
\end{tabular}}}\par
\wsec{Magnitude}
\Q{Find $|2\mathbf i-3\mathbf j+6\mathbf k|$.}
\vspace{2.6cm}
\Q{The points $A$ and $B$ have coordinates $(1,2,-1)$ and $(3,-4,2)$. Find the distance $AB$.}
\vspace{2.6cm}
\Q{\textbf{(Tricky)} Given $\mathbf a=\col{t}{2}{t-1}$ and $|\mathbf a|=3$, find the possible values of $t$.}
\vspace{2.6cm}
\par\needspace{8\baselineskip}\medskip\noindent\textbf{Check Your Understanding.} Circle T (True) or F (False) for each statement.\par\smallskip{\renewcommand{\arraystretch}{1.5}\noindent\makebox[\textwidth][c]{\begin{tabular}{|c|>{\centering\arraybackslash}p{12cm}|c|}\hline
 & \textbf{Statement} & \textbf{T/F}\\\hline
(a) & $|\mathbf i+2\mathbf j-2\mathbf k|=1+2-2=1$ & T\,/\,F \\[1.2ex]\hline
(b) & $|-\mathbf a|=|\mathbf a|$ for every vector $\mathbf a$ & T\,/\,F \\[1.2ex]\hline
(c) & $|\mathbf a+\mathbf b|=|\mathbf a|+|\mathbf b|$ for all vectors $\mathbf a$, $\mathbf b$ & T\,/\,F \\[1.2ex]\hline
(d) & $|3\mathbf a|=3|\mathbf a|$ & T\,/\,F \\[1.2ex]\hline
(e) & $\left|\col{0}{-5}{0}\right|=-5$ & T\,/\,F \\[1.2ex]\hline
(f) & $\left|\col{1}{1}{1}\right|=1$ & T\,/\,F \\[1.2ex]\hline
(g) & $|\mathbf a+\mathbf b|\le|\mathbf a|+|\mathbf b|$ for all vectors $\mathbf a$, $\mathbf b$ & T\,/\,F \\[1.2ex]\hline
\end{tabular}}}\par
\wsec{Unit Vectors}
\Q{Find the unit vector in the direction of $\col{2}{-1}{2}$.}
\vspace{2.6cm}
\Q{Find the vector of magnitude $10$ in the direction of $\col{3}{0}{-4}$.}
\vspace{2.6cm}
\Q{\textbf{(Tricky)} The points $A$ and $B$ have coordinates $(1,0,2)$ and $(3,-6,5)$. Find the vector of magnitude $14$ in the direction \emph{opposite} to $\overrightarrow{AB}$.}
\vspace{2.6cm}
\par\needspace{8\baselineskip}\medskip\noindent\textbf{Check Your Understanding.} Circle T (True) or F (False) for each statement.\par\smallskip{\renewcommand{\arraystretch}{1.5}\noindent\makebox[\textwidth][c]{\begin{tabular}{|c|>{\centering\arraybackslash}p{12cm}|c|}\hline
 & \textbf{Statement} & \textbf{T/F}\\\hline
(a) & $\col{1}{1}{1}$ is a unit vector & T\,/\,F \\[1.2ex]\hline
(b) & $\dfrac13\col{1}{2}{2}$ is a unit vector & T\,/\,F \\[1.2ex]\hline
(c) & $\mathbf i+\mathbf j$ is a unit vector & T\,/\,F \\[1.2ex]\hline
(d) & $|\hat{\mathbf a}|=1$ for every non-zero vector $\mathbf a$ & T\,/\,F \\[1.2ex]\hline
(e) & $\mathbf a$ and $\hat{\mathbf a}$ are parallel & T\,/\,F \\[1.2ex]\hline
\end{tabular}}}\par
\wsec{Parallel Vectors ($\mathbf a=\lambda\mathbf b$)}
\Q{Show that $\mathbf a=\col{2}{-1}{3}$ and $\mathbf b=\col{-6}{3}{-9}$ are parallel.}
\vspace{2.6cm}
\Q{Find $p$ and $q$ such that $\col{4}{p}{-2}$ is parallel to $\col{-2}{3}{q}$.}
\vspace{2.6cm}
\Q{\textbf{(Tricky)} Find all vectors of magnitude $6$ that are parallel to $\mathbf a=\col{1}{2}{-2}$.}
\vspace{2.6cm}
\par\needspace{8\baselineskip}\medskip\noindent\textbf{Check Your Understanding.} Circle T (True) or F (False) for each statement.\par\smallskip{\renewcommand{\arraystretch}{1.5}\noindent\makebox[\textwidth][c]{\begin{tabular}{|c|>{\centering\arraybackslash}p{12cm}|c|}\hline
 & \textbf{Statement} & \textbf{T/F}\\\hline
(a) & If $\mathbf a\parallel\mathbf b$, then $\mathbf a=\mathbf b$ & T\,/\,F \\[1.2ex]\hline
(b) & $\col{1}{-2}{3}$ is parallel to $\col{-2}{4}{-6}$ & T\,/\,F \\[1.2ex]\hline
(c) & $\col{2}{4}{6}$ is parallel to $\col{1}{2}{4}$ & T\,/\,F \\[1.2ex]\hline
(d) & $\col{3}{0}{-6}$ is parallel to $\col{-1}{0}{2}$ & T\,/\,F \\[1.2ex]\hline
(e) & If $\mathbf a=-2\mathbf b$, then $\mathbf a$ and $\mathbf b$ are not parallel since they point in opposite directions & T\,/\,F \\[1.2ex]\hline
\end{tabular}}}\par
\wsec{Collinear Points}
\Q{Show that the points $A(1,2,3)$, $B(3,5,4)$ and $C(7,11,6)$ are collinear.}
\vspace{2.6cm}
\Q{For the points in Question 1, find the ratio $AB:BC$.}
\vspace{2.6cm}
\Q{\textbf{(Tricky)} The points $P(2,-1,0)$, $Q(4,1,3)$ and $R(a,7,b)$ are collinear. Find $a$ and $b$.}
\vspace{2.6cm}
\par\needspace{8\baselineskip}\medskip\noindent\textbf{Check Your Understanding.} Circle T (True) or F (False) for each statement.\par\smallskip{\renewcommand{\arraystretch}{1.5}\noindent\makebox[\textwidth][c]{\begin{tabular}{|c|>{\centering\arraybackslash}p{12cm}|c|}\hline
 & \textbf{Statement} & \textbf{T/F}\\\hline
(a) & If $\overrightarrow{AB}\parallel\overrightarrow{CD}$, then $A$, $B$, $C$ and $D$ are collinear & T\,/\,F \\[1.2ex]\hline
(b) & If $\overrightarrow{AB}=3\overrightarrow{AC}$, then $A$, $B$ and $C$ are collinear & T\,/\,F \\[1.2ex]\hline
(c) & $A(0,0,0)$, $B(1,1,1)$ and $C(2,2,3)$ are collinear & T\,/\,F \\[1.2ex]\hline
(d) & The points with position vectors $\mathbf a$, $2\mathbf a$ and $-\mathbf a$ are collinear ($\mathbf a\neq\mathbf 0$) & T\,/\,F \\[1.2ex]\hline
\end{tabular}}}\par
\wsec{Ratio Theorem}
\begin{center}\begin{tikzpicture}[scale=0.85,>=Stealth]
\coordinate (O) at (0,0);\coordinate (A) at (0.6,2.4);\coordinate (B) at (5.2,1.6);
\coordinate (C) at ($(A)!0.4!(B)$);
\draw[thick,black!70](A)--(B);
\draw[va](O)--(A) node[midway,left]{\small$\mathbf a$};\draw[vb](O)--(B) node[midway,below]{\small$\mathbf b$};\draw[vc](O)--(C);
\node[dot,label=below left:$O$] at (O){};\node[dot,label=above left:$A$] at (A){};
\node[dot,label=right:$B$] at (B){};\node[dot,label=below:$C$] at (C){};
\draw[decorate,decoration={brace,amplitude=4pt,raise=3pt}](A)--(C) node[midway,above=8pt,sloped]{$\lambda$};
\draw[decorate,decoration={brace,amplitude=4pt,raise=3pt}](C)--(B) node[midway,above=8pt,sloped]{$\mu$};
\end{tikzpicture}\end{center}\noindent In the True/False table, $C$ lies on $AB$ with $AC:CB=\lambda:\mu$, $\overrightarrow{OA}=\mathbf a$ and $\overrightarrow{OB}=\mathbf b$.
\Q{The points $A$ and $B$ have position vectors $\col{1}{4}{-2}$ and $\col{7}{-2}{4}$. The point $C$ lies on $AB$ such that $AC:CB=1:2$. Find $\overrightarrow{OC}$.}
\vspace{2.6cm}
\Q{Find the position vector of the midpoint $M$ of $AB$, where $A$ and $B$ have coordinates $(3,-1,5)$ and $(-1,5,1)$.}
\vspace{2.6cm}
\Q{\textbf{(Tricky)} The point $A$ has coordinates $(1,2,3)$. The point $C(3,0,7)$ lies on $AB$ such that $AC:CB=2:3$. Find the coordinates of $B$.}
\vspace{2.6cm}
\par\needspace{8\baselineskip}\medskip\noindent\textbf{Check Your Understanding.} Circle T (True) or F (False) for each statement.\par\smallskip{\renewcommand{\arraystretch}{1.5}\noindent\makebox[\textwidth][c]{\begin{tabular}{|c|>{\centering\arraybackslash}p{12cm}|c|}\hline
 & \textbf{Statement} & \textbf{T/F}\\\hline
(a) & $\overrightarrow{OC}=\dfrac{\mu\mathbf a+\lambda\mathbf b}{\lambda+\mu}$ & T\,/\,F \\[1.2ex]\hline
(b) & $\overrightarrow{OC}=\dfrac{\lambda\mathbf a+\mu\mathbf b}{\lambda+\mu}$ & T\,/\,F \\[1.2ex]\hline
(c) & If $\lambda:\mu=1:1$, then $\overrightarrow{OC}=\dfrac12(\mathbf a+\mathbf b)$ & T\,/\,F \\[1.2ex]\hline
(d) & If $\lambda:\mu=1:3$, then $\overrightarrow{OC}=\dfrac{3\mathbf a+\mathbf b}{4}$ & T\,/\,F \\[1.2ex]\hline
(e) & If $\lambda:\mu=2:1$, then $\overrightarrow{OC}=\dfrac{2\mathbf a+\mathbf b}{3}$ & T\,/\,F \\[1.2ex]\hline
\end{tabular}}}\par
\par\vspace{\baselineskip}\begin{center}\textbf{END OF WORKSHEET}\end{center}\end{document}
